\documentclass[10pt]{article}
\usepackage[a4paper,margin=22mm]{geometry}
\usepackage{amsmath,amssymb,amsthm,mathtools}
\usepackage{booktabs}
\usepackage{microtype}
\usepackage{enumitem}
\usepackage{xurl}
\usepackage[hidelinks]{hyperref}
\providecommand{\CRevisionRound}{2}
\pdfvariable suppressoptionalinfo 767
\setlength{\parindent}{1.2em}
\setlength{\parskip}{0.15em}
\setlist{nosep,leftmargin=1.8em}
\hbadness=10000
\vbadness=10000
\emergencystretch=2em

\newtheorem{theorem}{Theorem}
\newtheorem{lemma}{Lemma}
\newtheorem{proposition}{Proposition}
\theoremstyle{remark}
\newtheorem{remark}{Remark}
\newcommand{\R}{\mathcal R}
\newcommand{\e}{\mathrm e}

\title{One Lyapunov Ratio, Three Global Phases:\
The May--Leonard Coexistence--Cycle Atlas}
\author{HCS-C358 source-local reconstruction}
\date{3 September 2026}

\begin{document}
\maketitle

\begin{abstract}
For the three-species cyclic May--Leonard competition flow in the strict
intransitive chamber, we prove an all-interior-orbit classification.  The
normalized product $\R=xyz/(x+y+z)^3$ has a derivative whose sign is exactly
that of $2-a-b$.  Below the critical plane it yields global convergence to
coexistence.  On the plane it combines with an exact logistic time change to
give a foliation by periodic rock--paper--scissors leaves and an elliptic
period quadrature.  Above the plane, the positive diagonal is the complete
stable set of the unstable coexistence point, while every other interior
orbit has the full oriented boundary heteroclinic cycle as its omega-limit
set and has diverging residence times.  Exact rational and symbolic receipts
audit the identities and conventions without replacing the global proof.  No
literature-priority or arithmetic-target claim is made.
\end{abstract}

\noindent\textbf{Revision certificate.}
\ifnum\CRevisionRound=0
Round-zero invariant and coexistence closure.
\else\ifnum\CRevisionRound=1
Round-one critical foliation and elliptic-period closure.
\else
Round-two heteroclinic exhaustion, evidence, and Route-A closure.
\fi\fi

\small
\section{Model, chamber, and main result}

For $a,b\geq0$, consider
\begin{equation}
\begin{aligned}
 \dot x&=x(1-x-ay-bz),\\
 \dot y&=y(1-bx-y-az),\\
 \dot z&=z(1-ax-by-z).
\end{aligned}
\end{equation}
We work in the strict cyclic chamber
\begin{equation}
                 (a-1)(b-1)<0.
\end{equation}
Put $S=x+y+z$, $Q=xy+yz+zx$, and, in the interior,
$\R=xyz/S^3$.  Let
\[
 p=\frac{1}{1+a+b}(1,1,1),\qquad
 E_1=(1,0,0),\ E_2=(0,1,0),\ E_3=(0,0,1).
\]

\begin{theorem}[Complete cyclic-chamber atlas]\label{thm:atlas}
The closed positive octant is forward invariant; its interior is preserved;
and every positive semiorbit is bounded and forward complete.  The point
$p$ is the unique interior equilibrium.  For every interior initial state:
\begin{enumerate}[label=\textup{(\roman*)}]
\item If $a+b<2$, the solution converges to $p$.
\item If $a+b=2$, then $S$ obeys a logistic equation.  With
$(u,v,w)=(x,y,z)/S$, $\delta=1-a=b-1$, and $d\tau=S\,dt$, the normalized
flow is
\begin{equation}
 u_\tau=\delta u(v-w),\quad
 v_\tau=\delta v(w-u),\quad
 w_\tau=\delta w(u-v).
\end{equation}
Its centre is fixed, and each level $uvw=h$, $0<h<1/27$, is one periodic
orbit.  If $r_-(h)<1/3<r_+(h)$ are the roots in $(0,1)$ of
$r(1-r)^2=4h$, its least $\tau$-period is
\begin{equation}
 T(h)=\frac{2}{|\delta|}\int_{r_-(h)}^{r_+(h)}
 \frac{dr}{\sqrt{r\{r(1-r)^2-4h\}}}.
\end{equation}
Every original noncentral solution approaches a unique phase translate of
its normalized periodic orbit.
\item If $a+b>2$, the complete interior stable set of $p$ is the positive
diagonal.  Every other interior orbit has as its omega-limit set the full
three-connection boundary heteroclinic cycle and spends unbounded successive
times near its axial saddles.  For $a<1<b$ its orientation is
$E_1\to E_3\to E_2\to E_1$; for $b<1<a$ the arrows reverse.
\end{enumerate}
At the closure point $a=b=1$, all ratios are fixed and the simplex $S=1$
consists entirely of equilibria.  The separate walls $a=1$ or $b=1$ are
nonhyperbolic and are not included in (2).
\end{theorem}

\section{Positive flow and the source identity}

Each coordinate in (1) is itself times a locally Lipschitz function.  Thus a
zero coordinate remains zero, whereas a positive one cannot vanish at finite
time.  Since $a,b\geq0$,
\[
 \dot x\leq x(1-x),\qquad \dot y\leq y(1-y),\qquad
 \dot z\leq z(1-z).
\]
Scalar comparison proves boundedness and continuation for all positive time.
For some $C=C(a,b)>0$, summing (1) also gives
$\dot S\geq S-CS^2$.  Consequently an interior semiorbit eventually lies in
a compact annulus bounded away from the origin.

At an interior equilibrium,
\[
 \begin{pmatrix}1&a&b\\ b&1&a\\ a&b&1\end{pmatrix}
 \begin{pmatrix}x\\y\\z\end{pmatrix}
 =\begin{pmatrix}1\\1\\1\end{pmatrix}.
\]
The determinant is
\[
 (1+a+b)(1+a^2+b^2-a-b-ab),
\]
and twice the second factor is
$(a-b)^2+(a-1)^2+(b-1)^2$.  It vanishes only at $a=b=1$.
Subtracting cyclic rows therefore gives the sole positive solution $p$ in
the strict chamber.

\begin{lemma}[Normalized-product identity]\label{lem:key}
Along every interior solution,
\begin{align}
 \dot S&=S-S^2+(2-a-b)Q,\\
 \frac{d}{dt}\log(xyz)&=3-(1+a+b)S,\\
 \frac{d}{dt}\log\R&=(2-a-b)\frac{S^2-3Q}{S},\\
 S^2-3Q&=\tfrac12\{(x-y)^2+(y-z)^2+(z-x)^2\}.
\end{align}
\end{lemma}

\begin{proof}
Adding (1) and using $x^2+y^2+z^2=S^2-2Q$ gives (5).  Dividing each equation
by its coordinate and adding gives (6).  Equation (7) is (6) minus
$3\dot S/S$, and direct expansion gives (8).
\end{proof}

The derivative in (7) is strict away from the positive diagonal whenever
$a+b\ne2$.  In particular there can be no nonconstant interior periodic
orbit off the critical plane.

\section{Subcritical global coexistence}

Assume $a+b<2$.  By Lemma~\ref{lem:key}, $\R$ increases and is bounded above
by $1/27$.  The bound $\R(t)\geq\R(0)>0$ keeps $(x/S,y/S,z/S)$ away from the
simplex boundary, while the annulus estimate controls $S$.  Hence the orbit
is eventually contained in a compact subset of the interior.  LaSalle's
invariance principle applied to $\log\R$ places its omega-limit set on
$x=y=z$.  That diagonal is invariant, and $x=y=z=r$ reduces (1) to
\begin{equation}
             \dot r=r\{1-(1+a+b)r\}.
\end{equation}
The only compact complete diagonal orbit in the positive annulus is the
equilibrium $r=(1+a+b)^{-1}$.  Thus the omega-limit set is $\{p\}$ and the
whole orbit converges to $p$, proving Theorem~\ref{thm:atlas}(i).

\ifnum\CRevisionRound>0
\section{Critical normalization and the elliptic period}

Now let $a+b=2$.  Equation (5) becomes $\dot S=S(1-S)$, so, with $S_0=S(0)$,
\begin{equation}
 S(t)=\frac{S_0\e^t}{1+S_0(\e^t-1)},\qquad
 \tau(t)=\log\{1+S_0(\e^t-1)\}.
\end{equation}
Direct differentiation gives $\dot\tau=S$ and
$\tau(t)-t\to\log S_0$.

Differentiating $u=x/S$, cyclically, proves (3).  The components of (3) sum
to zero and $d(uvw)/d\tau=0$.  On the open simplex, the sole equilibrium is
the centre.  For $0<h<1/27$, the regular set
\[
 \Gamma_h=\{u+v+w=1: u,v,w>0,\ uvw=h\}
\]
is a compact circle surrounding the centre.  Its vector field never
vanishes, so it is exactly one periodic orbit.

To compute the period, use $v+w=1-u$ and $vw=h/u$:
\begin{equation}
 (u_\tau)^2=\delta^2u^2(v-w)^2
 =\delta^2u\{u(1-u)^2-4h\}.
\end{equation}
The function $r(1-r)^2$ rises from zero to $4/27$ on $(0,1/3)$ and falls
back to zero on $(1/3,1)$.  Its two roots at height $4h$ are precisely the
turning values $r_-,r_+$.  Integrating (11) from one turning point to the
other and back yields (4); this is the least period because the vector field
traverses $\Gamma_h$ once.

Let $U(\tau)$ be the periodic normalized solution with
$U(0)=(u(0),v(0),w(0))$.  Equation (10) gives
$(x,y,z)(t)=S(t)U(\tau(t))$.  Since $S(t)\to1$ and
$\tau(t)-(t+\log S_0)\to0$, uniform continuity of $U$ gives
\begin{equation}
 (x,y,z)(t)-U(t+\log S_0)\longrightarrow0.
\end{equation}
This proves the stated asymptotic phase, uniquely modulo the least period.
\fi

\ifnum\CRevisionRound>1
\section{The diagonal exception above criticality}

Let $a+b>2$.  Linearization at $p$ has the radial eigenvalue $-1$ and a
complex tangent pair satisfying
\begin{equation}
 \operatorname{Re}\lambda=\frac{a+b-2}{2(1+a+b)}>0,
 \qquad
 (\operatorname{Im}\lambda)^2=\frac{3(a-b)^2}{4(1+a+b)^2}.
\end{equation}
Thus the local stable manifold is one-dimensional.  The positive diagonal
is invariant and tangent to the stable eigenspace.  Uniqueness in the stable
manifold theorem identifies them locally.  If an interior orbit converges to
$p$, it eventually belongs to this local stable manifold; uniqueness of the
flow then puts its entire orbit on the diagonal.  Conversely, (9) shows that
every positive diagonal orbit converges to $p$.  Hence the complete interior
stable set is exactly the positive diagonal.

For every other interior orbit, (7) makes $\R$ strictly decreasing.  If its
limit were positive, the normalized coordinates would remain away from the
boundary.  LaSalle's principle applied to $-\log\R$ would then force the
omega-limit set to $p$, contradicting the stable-set classification.  Thus
\begin{equation}
                         \R(t)\longrightarrow0.
\end{equation}
Since $S$ stays in a positive compact annulus, the omega-limit set is a
nonempty compact invariant subset of the nonzero boundary.

\section{Boundary exhaustion and residence times}

Suppose first that $a<1<b$.  On $z=0$,
\begin{equation}
 \frac{d}{dt}\log\frac{y}{x}=-(b-1)x-(1-a)y<0.
\end{equation}
There is no positive equilibrium in this coordinate plane.  The Dulac
multiplier $(xy)^{-1}$ has divergence $-1/x-1/y<0$ there, so there is no
positive planar cycle.  By Poincar\'e--Bendixson, the absence of a positive
planar equilibrium, and the strict ratio in (15), the only nonconstant
bounded complete planar orbit whose closure avoids the origin is the positive
branch of the unstable manifold of $E_2$, connecting $E_2$ to $E_1$.
Cyclic permutation gives the three connections
\[
 E_1\longrightarrow E_3\longrightarrow E_2\longrightarrow E_1.
\]
Every other positive-plane orbit either has the origin in its backward
closure or escapes compact sets in negative time.  A non-equilibrium axial
orbit also has the origin in its backward closure.

Let $\Omega$ be the omega-limit set of the interior orbit.  Standard flow
properties make $\Omega$ compact, connected, invariant, and internally chain
transitive.  It is bounded away from the origin.  The preceding planar
classification and (14) therefore confine $\Omega$ to the three axial
equilibria and the three connections (16).  It cannot be one axial point:
at each $E_i$ the stable manifold in the closed octant is the incoming
coordinate face, whereas the omitted-species eigenvalue is $1-a>0$.  Nor can
it be a proper closed directed subchain.  Choose disjoint small
neighbourhoods of the axial points and compact middle arcs on the included
connections.  On each middle arc the derivative in (15), or its cyclic
analogue, has a uniform strict sign, and the flow crosses the two boundary
sections only in the directed order.  For sufficiently small jump size, no
long-time pseudo-orbit within a proper subchain can return from its terminal
axial neighbourhood to its initial one.  Internal chain transitivity leaves
exactly the full cycle (16).  This proves
$\Omega$ equals that cycle, not merely that its distance from the boundary
tends to zero.

The orbit enters arbitrarily small fixed neighbourhoods of every $E_i$.
At successive entries its outgoing coordinate tends to zero.  In such a
neighbourhood that coordinate grows no faster than
$1-a+\varepsilon$.  The time needed to reach a fixed exit section is thus at
least a constant multiple of the logarithm of the reciprocal entry
coordinate, and tends to infinity.  Swapping $y$ and $z$ swaps $a,b$ and
reverses (16), completing Theorem~\ref{thm:atlas}(iii).

\section{Closure faces, evidence, and claim boundary}

On each coordinate axis, (1) is the scalar logistic flow.  The origin is
repelling relative to the positive octant.  At $a=b=1$, all equations have
the form $\dot x_i=x_i(1-S)$: ratios are constant, $S$ is logistic, and the
whole simplex $S=1$ is equilibria.  On either separate wall $a=1$ or $b=1$,
an axial transverse eigenvalue vanishes; those dominance walls are recorded
but excluded from the strict theorem.

The exact certificate contains 12 parameter rows, 72 invariant rows, 24
critical rows, 22 logistic rows, six orientation rows, and five boundary
rows.  An independent implementation reconstructs all rows and strict
serialization contracts; a separate symbolic derivation checks 19
identities; two isolated replays agree byte for byte; and 67 repaired-hash
mutations are rejected.  These are receipts for algebra and convention.
They do not prove LaSalle exhaustion, stable-manifold uniqueness, periodic
foliation, or the omega-limit theorem; the preceding arguments do.

May and Leonard supply the original cyclic-competition lineage
\cite{mayleonard}; Bl\'e, Castellanos, Llibre, and Quilant\'an provide later
integrability and global dynamics context \cite{bleetal}.  We make no
priority claim.

The strict Route-A tuple is
\[
 (\mathrm{A0\_FAIL},\mathrm{A1\_WEAK},\mathrm{A2\_FAIL},
  \mathrm{A3\_FAIL},\mathrm{A4\_FAIL}).
\]
The source has a continuum of critical periodic leaves but no intrinsic
rational-prime objects, prime-power repetition law, logarithmic-prime clock,
target determinant, or natural target-zero quantization.  Route A is rejected
and Route B is locked.  In particular, no target arithmetic local data,
Euler factors, bad-prime data, root number, automorphy, target divisor or
counting law, target functional equation, target-zero match, or
Hilbert--P\'olya operator is claimed.

\begin{thebibliography}{9}
\bibitem{mayleonard}
R. M. May and W. J. Leonard,
``Nonlinear aspects of competition between three species,''
\emph{SIAM J. Appl. Math.} 29 (1975), 243--253.
\href{https://doi.org/10.1137/0129022}{doi:10.1137/0129022}.

\bibitem{bleetal}
G. Bl\'e, V. Castellanos, J. Llibre, and I. Quilant\'an,
``Integrability and global dynamics of the May--Leonard model,''
\emph{Nonlinear Anal. Real World Appl.} 14 (2013), 280--293.
\href{https://doi.org/10.1016/j.nonrwa.2012.06.004}{doi:10.1016/j.nonrwa.2012.06.004}.
\end{thebibliography}
\fi

\end{document}
